\subsection{محدودیت‌ها و تحلیل شکست}
تحلیل مکانیزم‌های خطی‌سازی ماژول \lr{liftr} محدودیت‌های بنیادی زیر را نمایان می‌سازد:
\begin{itemize}
    \item \textbf{فقدان فشرده‌سازی فیزیکی پارامترها (\lr{No Parameter Reduction}):} ماژول \lr{liftr} ماهیت فشرده‌سازی محاسباتی (\lr{Compute Complexity Reduction}) دارد و ماتریس‌های وزن را حذف یا هرس نمی‌کند؛ بنابراین حافظه ایستا اشغال‌شده توسط وزن‌های مدل بدون تغییر باقی می‌ماند.
    \item \textbf{از بین رفتن وابستگی‌های مکانی توکن‌ها (\lr{Loss of Spatial Resolution}):} تبدیل ماتریس توجه به میانگین یکنواخت مقادیر ارزش‌ها موجب می‌شود که اطلاعات مربوط به ترتیب و وزن‌دهی اختصاصی میان جفت توکن‌ها از بین برود؛ این امر در تسک‌هایی که نیازمند استدلال مکانی یا ردیابی دقیق مؤلفه‌ها در بافت‌های پیچیده هستند موجب افت کیفیت بازنمایی می‌گردد.
    \item \textbf{ناسازگاری با سرهای توجه گروهی در مدل \lr{Qwen} (\lr{GQA Incompatibility}):} عدم تطابق نسبت تعداد سرهای کوئری به کلید/ارزش در ساختار \lr{Grouped-Query Attention} منجر به بروز خطا در عملگر پخش (\lr{Broadcast}) ابعاد میانگین ارزش‌ها شده که عامل ریشه‌ای شکست ارزیابی در مدل \lr{Qwen} بوده است.
\end{itemize}
